
ML benchmarks lose their ability to discriminate between methods once the community has converged on high-performing approaches --- yet decisions about when to retire them remain informal and subjective. This paper shows that benchmark saturation is not a gradual decline but a discrete structural break. Among 115 Papers With Code benchmarks, PELT change-point detection on OLS-detrended residual coefficient of variation (CoV) identifies a statistically significant regime shift at paper\_count\* $\approx$ 39 (permutation p=0.035, piecewise F-test p=0.0021). Before this threshold, performance variance is nearly $4\times$ higher than the global baseline; after it, post-breakpoint variance collapses to 20\% of pre-breakpoint levels (an 80\% reduction relative to the pre-breakpoint segment) --- consistent with Goodhart saturation dynamics, where community convergence on dominant approaches compresses the score distribution toward a ceiling. Three experiments with distinct test designs and null hypotheses corroborate this two-regime structure (pre-regime F p=0.0009, post-regime Brown-Forsythe p=0.0099). The results provide the first data-driven, paper-count-based structural break threshold for PwC benchmarks, offering a lightweight, reproducible tool that any leaderboard curator can run on existing data to flag mature benchmarks for retirement review.

